\section{Discussion and Related Work}
\subsection{What the Evidence Establishes}
The full factorial answers which measured configuration wins on each Amazon
workload; it does not supply a query-independent oracle choosing those winners.
Changing one upper topology at fixed base and entry provider is controlled.
The 0L matched-provider comparison instead measures an integrated topology and
entry-design choice. Its favorable low-selectivity results do not extend to
all broad predicates. Fewer logical edges can reduce navigation alternatives,
while a larger prefix frontier increases seed work. Additional overlays
introduce authorization and seed competition even with isolated adjacency.

The automatic policy provides a reproducible starting configuration without
query-performance calibration. It does not automatically select the base,
and its mass gate is a structural heuristic rather than a latency model.
All four evaluated datasets derive two overlays, so threshold/topology transfer
is tested but depth-change performance is not. The VariousImg regression is
retained. A submission-strength evaluation still needs current-protocol
external baselines, longer repeated timing, and query-quality checks for each
accelerated construction profile. Historical external runs use a different
protocol and are not merged with the factorial. Predicate safety here concerns
static AND-containment indexes with root-prefix authorization; extending the
predicate language or update model requires a new contract.

\subsection{Filtered and Range-Constrained Graph Search}
HNSW~\cite{malkov2020hnsw} and DiskANN~\cite{subramanya2019diskann} organize
vector proximity for bounded ANN search. Filtered-DiskANN
~\cite{gollapudi2023filtereddiskann} incorporates labels during graph
construction. ACORN~\cite{patel2024acorn} supports predicate-agnostic graph
search with expanded neighborhoods and compressed construction; its design
includes a configurable low-selectivity pre-filtering fallback. Its operating regime therefore depends on both the graph and the fallback
configuration.

UNG~\cite{cai2024ung} is the immediate structural baseline. Its exact-label
groups, minimal-superset LNG, and vector realizations jointly preserve
containment-directed navigation. We expose a different entry/topology contract:
a terminal-prefix forest with its covering frontier. The forest has fewer
permitted relations than general containment, but its edges are not necessarily
an edge-subset of the LNG cover graph. We then aggregate compatible groups
into independently authorized blocks. HNSW instead samples vector vertices
across geometric layers; our scales change the label-space navigation unit.

SeRF~\cite{zuo2024serf}, iRangeGraph~\cite{xu2024irangegraph}, and
UNIFY~\cite{liang2024unify} compress, compose, or combine interval-specialized
graphs. Dynamic range systems extend such structures to arrival and update
workloads~\cite{peng2025dynamic,zhang2025rangepq}. Their ordered scalar
intervals inspire explicit granularity, but arbitrary label containment forms
a partial order and requires a different authorization argument.

Curator~\cite{jin2026curator} shares label-specialized virtual indexes through
a clustering tree. Its temporary indexes for conjunctions can be assembled by
qualified-ID ordering and partition lookup without distance computation; its
reported two-label AND results are close to pre-indexed predicates. Our
blocks instead have fixed canonical prefix membership and use a root-prefix
certificate to authorize traversal. FAVOR~\cite{song2026favor} selects a filter-first
brute-force path below an estimated selectivity threshold and an
exclusion-distance HNSW path otherwise. Its empirically chosen threshold and
reported dataset-dependent crossover motivate studying multiple regimes;
the absolute candidate mass remains relevant to the cost of its scan route.

Elastic Index Selection~\cite{yang2025elastic} and SIEVE~\cite{li2025sieve}
choose and share collections of indexes. Their selection problem differs from
choosing block scales inside one containment-oriented system. Recent
LSSG~\cite{wang2026lssg}, an arXiv preprint, stratifies edges by label-set
similarity and traverses multiple tiers. Thus merely adding label-aware layers
is not our novelty: our distinctive mechanism is direct-member block
aggregation, exact root-prefix authorization, and states restricted to one
physical adjacency level. Query-aware Routing~\cite{xiong2026routing}, also
a preprint, uses an offline performance table and a learned Recall model;
DRH deliberately uses neither, accepting the resulting loss of adaptability.
These preprints are distinguished from accepted SIGMOD/PVLDB publications.

\subsection{Construction and Maintenance}
NN-Descent~\cite{dong2011nndescent} and relative-neighborhood pruning
~\cite{toussaint1980rng} provide building blocks for proximity graphs.
Database work improves construction with tree/graph hybrids, LSH-assisted
insertion, monotonic neighborhoods, and revised pruning
~\cite{azizi2023elpis,zhao2023lshapg,peng2023taumng,yang2025pgbuild}.
Starling~\cite{wang2024starling} co-designs navigation and storage for disk
I/O. GPU similarity search~\cite{johnson2021gpu} emphasizes distance and
selection throughput. Our contribution is to batch the irregular local-graph
and cross-edge tasks induced by a metadata hierarchy, retaining an explicit
host-output boundary. The static experiments do not establish dynamic-update
cost or numerical-quality equivalence across construction backends.

\section{Conclusion}
Entry discovery and graph topology form a joint contract in label-filtered
ANN. A prefix frontier covers all eligible terminal branches at the label
plane, while independently partitioned blocks admit coarser navigation when
a root-prefix certificate holds. Level-tagged states make the adjacency used
by each expansion explicit. These structural properties explain the design;
measured Recall and timing establish where it is useful. The Amazon factorial
shows both substantial improvements and topology-dependent regressions.
DRH supplies a query-calibration-free starting plan, and batched GPU
construction reduces measured build cost. Neither is claimed to dominate
manual tuning or preserve quality without its corresponding evidence.
